% ============================================================
% 5 实验
% ============================================================
\section{实验}
\label{sec:exp}

\subsection{实验设置}

数据与划分见第 3.2 节。除特别说明外，评测协议为图库 $1600$ 张测试集真实帧、
查询 $8000$ 张测试集伪造帧，报告帧级与视频级 Top-1$/$Top-5。
实现细节：图像尺寸 $256\times256$，描述子嵌入维度 $2048$，训练采用 Adam 优化器与
梯度检查点（Windows 环境下 \code{num\_workers=0}）。编码批次 $4\sim8$。
全部实验在同一测试协议下进行，随机基线帧级 Top-1 为 $0.0625\%$。
第 5.7 节另用一套稠密帧库协议（图库 $3000$ 帧 / 查询 $540$ 张），其余各节均为上述协议。

\subsection{参照方法与零训练基线}
\label{sec:baseline}

表~\ref{tab:hashvae} 列出主动检索框架 HashVAE 在其原文中报告的检索与检测结果。
我们曾尝试复现该框架，但未能达到原文所报告的水平（其训练依赖前景$/$背景的像素级重构、
解析式掩膜与更大规模的数据），因此本文不把复现读数当作该方法的性能，
统一以原文数值作为参照。

需要说明，原文的评测设定与本文并不一致：每个视频只随机抽取一帧，训练$/$测试按 $6{:}2{:}2$ 划分，
检索指标为 Top-5，且检索方向为由真实图查询伪造库；本文则是以伪造帧查询真实图库。
两组数值因此不能直接比较，只能作为量级参照。

\begin{table}[t]
\centering
\caption{HashVAE 原文报告的结果（数据取自原文表 2 与表 3，本文未复现该框架）。}
\label{tab:hashvae}
\small
\begin{tabular}{llcc}
\toprule
来源 & 设定 & 指标 & 数值 \\
\midrule
原文表 2 & FF++\_Deepfake, c23, 64 bit  & Top-5 & 84.83 \\
原文表 2 & FF++\_Deepfake, c23, 128 bit & Top-5 & 88.14 \\
原文表 2 & FF++\_Deepfake, c23, 1024 bit & Top-5 & 98.23 \\
原文表 3 & Deepfake (C40) 检测 & ACC / AUC & 96.25 / 96.62 \\
\bottomrule
\end{tabular}
\end{table}

表~\ref{tab:zero} 给出在本文协议下\textbf{完全不训练}的描述子读数。
零训练描述子无须拟合任何参数，因而可以直接回答“该任务本身有多难”。
其中 $8\times8$ 逐图标准化缩略图（$192$ 维，记为 \code{r8\_norm}）在帧级达到 $60.05\%$，
是本协议下最强的帧级基线；其余描述子（$32\times32$ 缩略图、颜色直方图、
冻结 DINOv2 特征）均在 $55\%\sim59\%$ 之间。

\begin{table}[t]
\centering
\caption{零训练描述子在统一协议（图库 1600 / 查询 8000）下的读数（\%）。}
\label{tab:zero}
\small
\begin{tabular}{llcccc}
\toprule
方法 & 区域 & 帧 Top-1 & 帧 Top-5 & 视频 Top-1 & 视频 Top-5 \\
\midrule
pHash（64 bit）              & 背景 & 49.34 & 78.42 & 94.26 & 97.61 \\
$32\times32$ 缩略图          & 背景 & 55.79 & 85.14 & 99.11 & 99.44 \\
$32\times32$ 缩略图          & 整图 & 58.80 & 87.15 & 99.75 & 99.80 \\
颜色直方图                   & 整图 & 58.23 & 87.33 & 99.90 & 99.90 \\
冻结 DINOv2 ViT-S/14\cite{oquab2024dinov2} & 整图 & 57.05 & 87.56 & 99.91 & 99.91 \\
冻结 DINOv2 ViT-S/14         & 背景 & 48.96 & 83.81 & 99.72 & 99.86 \\
\textbf{r8\_norm（192 维）}   & 整图 & \textbf{60.05} & -- & -- & -- \\
\bottomrule
\end{tabular}
\end{table}

\textbf{结论。}（i）该任务在帧级上并不依赖训练：一个 $192$ 维、逐图标准化的 $8\times8$
缩略图即可达到 $60.05\%$，远高于随机基线 $0.0625\%$，说明判别信号主要存在于
\textbf{低频背景}中。（ii）在视频级上所有零训练描述子都接近 $99\%$，
这是“图库与查询按帧序号对齐”的直接结果：同源帧之间几乎逐像素相同，
视频级指标已接近饱和，帧级才是真正有分辨力的层级。
（iii）这两点共同意味着，任何学习方法都应先与零训练基线对照，
才能说明其收益来自“学习”本身。

\subsection{帧级失效的成因：假设与证据}
\label{sec:rounds}

在同一协议下，我们训练了第 4.2 节的监督对比编码器，并对监督粒度、人脸区域擦除、
多尺度跳连、几何增强、低层蒸馏与嵌入头结构等做了十余轮迭代。
这些迭代暴露出一个稳定的现象：\textbf{视频级始终接近饱和，帧级却长期停在 $46\%\sim48\%$}；
最好的一轮（I-best）也只有 $59.68\%$，低于零训练基线 $60.05\%$。
容量不是解释——加深主干、延长训练、更换监督粒度都只改变几个百分点。

据此我们提出假设：帧级失效的原因是建模过程主动抹去了区分“同一视频内近乎相同帧”
所需的\textbf{微小空间信号}。四条读数支持这一假设：

\begin{enumerate}[leftmargin=1.6em,itemsep=1pt]
\item \textbf{全局池化。} 帧级瓶颈位于嵌入头之前的全局平均池化。去掉它、保留空间布局后，
      帧级由约 $46\%$ 升至约 $58\%$（$+12$ 个百分点），说明有用信息确实在空间布局里。
\item \textbf{几何增强。} 在训练视图中加入随机旋转、缩放与平移的一轮，
      首轮帧级由约 $40\%$ 降至约 $32\%$。几何增强要求模型对位置变化不变，
      而帧级恰恰要靠位置差异判别，两者方向相反。
\item \textbf{零训练空间网格。} 完全不训练、但保留 $8\times8$ 网格并逐图标准化的描述子
      达到 $60.05\%$，强于所有学习到的全局向量。说明“保留空间网格”比“学习”更关键。
\item \textbf{同视频 $8$ 选 $1$ 诊断。} 把候选集限制为查询自身源视频的那 $8$ 张真实帧
      （随机基线 $12.5\%$），即剥离视频级信息、只留“同视频挑对那一帧”：
      零训练 \code{r8\_norm} 为 $60.11\%$，学习编码器为 $59.79\%$，两者持平。
      而标准协议下学习编码器的帧级 Top-1 为 $59.68\%$，与 $8$ 选 $1$ 几乎相同，
      说明跨视频几乎不出错，帧级的全部损失都发生在同视频内部。
\end{enumerate}

第四条把“学习抹掉了细节”修正为更准确的表述：学习\textbf{并未破坏}细节，
而是只追平了零训练描述子，未能在此基础上增加分辨力。为了排除“负样本不够难”这一解释，
我们又做了一轮单变量训练：每批固定取同一条视频的两个不同源帧，使“背景逐像素相同、
只有人脸微动”的相邻帧必然同时出现在批内，其余超参与 I-best 逐项相同。
三轮读数为 $58.99\%$、$57.70\%$、$58.09\%$，$8$ 选 $1$ 口径下为 $59.10\%$，
均不优于 I-best（$59.68\%$ / $59.79\%$）。也就是说，\textbf{难度不是瓶颈}：
即便把最难的负样本直接送进批内，单一全局向量的比较方式仍然停在零训练水平。

第二个待检验的解释是\textbf{比较形式}本身。我们做了一轮单变量训练：把嵌入头之前的池化张量
按 $8\times8$ 网格拆成 $64$ 个 $32$ 维局部描述子，在\textbf{同一空间位置}上做跨样本监督对比
（温度 $0.1$、权重 $0.5$），其余超参与 I-best 逐项相同，使网络被直接要求“每个格子都具备
区分近重复帧的能力”。训练期内该损失由 $0.55$ 降至 $0.11\sim0.29$，说明逐位置区分确实在学；
但三轮帧级读数为 $60.05\%$、$59.24\%$、$59.55\%$，$8$ 选 $1$ 口径为 $60.14\%$。
它相比 I-best（$59.68\%$ / $59.79\%$）只高约 $0.35$ 个百分点——$8000$ 张查询中约 $28$ 张，
小于同配置不同轮次之间约 $0.8$ 个百分点的波动，也与零训练 \code{r8\_norm} 的
$60.05\%$ / $60.11\%$ 持平。需要强调的是，这一轮的检索仍然使用 head 混合后的单一全局向量，
逐位置监督只作为辅助损失，因此它检验的是“给全局向量补一条位置级监督”，
而\textbf{不是}“改变检索时的比较方式”；读数说明前者无效，并不否定后者。

同一批零训练扫描也解释了为什么更细的网格不会自动带来更多信息：逐图标准化缩略图的帧级命中
在 $8\times8$ 处达到峰值，此后随分辨率上升反而下降（$4\times4$ 为 $58.61\%$、
$8\times8$ 为 $60.05\%$、$16\times16$ 为 $59.35\%$、$32\times32$ 为 $58.70\%$、
$64\times64$ 为 $58.06\%$），而只保留高频的输入仅 $46.74\%$。
也就是说，判别信息集中在中低分辨率的空间布局上，靠提高分辨率或提取高频都拿不到。

真正改变的是检索时的比较方式。第 5.6 节用同一批候选做逐位置的软匹配与几何校验，
同一个 $8$ 选 $1$ 口径下命中率由 $59.79\%$ 升至 $66.25\%$，干净档标准协议下由
$60.05\%$ 升至 $67.67\%$。该提升中“更细的局部网格”“显式对应”与“局部描述子的独立训练”
各自的贡献本文尚未分离，留作下一步。

\subsection{退化鲁棒基准}
\label{sec:degrade-exp}

表~\ref{tab:degrade} 给出 12 档退化下闭集子集（图库 $1200$ 张真实帧，对应第 4.5 节开集协议的
闭集部分）的帧级读数，被测对象为学习编码器 I-best 与零训练 \code{r8\_norm}。

\begin{table}[t]
\centering
\caption{12 档确定性退化下的帧级读数（只退化查询；图库干净，$1200$ 帧）。}
\label{tab:degrade}
\small
\begin{tabular}{lcccc}
\toprule
\multirow{2}{*}{档位} & \multicolumn{2}{c}{I-best} & \multicolumn{2}{c}{r8\_norm（零训练）} \\
\cmidrule(lr){2-3}\cmidrule(lr){4-5}
 & 帧 Top-1 & 帧 Top-5 & 帧 Top-1 & 帧 Top-5 \\
\midrule
clean      & 60.36 & 89.10 & 61.03 & 88.90 \\
jpeg50     & 60.23 & 89.04 & 61.16 & 88.77 \\
jpeg20     & 60.36 & 88.57 & 61.16 & 88.84 \\
scale0.5   & 60.29 & 88.97 & 61.03 & 88.84 \\
crop0.5    & 1.34  & 3.81  & 0.47  & 2.41 \\
reencode   & 60.29 & 88.97 & 61.03 & 88.57 \\
\midrule
crop0.8    & 23.86 & 69.92 & 20.59 & 62.10 \\
crop0.7off & 8.82  & 34.29 & 4.48  & 18.78 \\
zoom1.15   & 34.22 & 81.75 & 33.16 & 80.35 \\
rot5       & 48.33 & 87.10 & 44.92 & 87.30 \\
rot10      & 34.56 & 83.29 & 30.88 & 80.95 \\
shift8     & 38.03 & 83.36 & 36.83 & 83.09 \\
\bottomrule
\end{tabular}
\end{table}

\textbf{结论。}（i）\textbf{光度六档几乎不改变读数}：clean 与 jpeg50/jpeg20/scale0.5/reencode
的帧 Top-1 差异均在 $0.2$ 个百分点以内，这印证了第 5.2 节的判断——判别信号在低频，
而光度退化只削弱高频细节。（ii）\textbf{几何六档才是有效区分者}：旋转 $5^\circ$ 已使帧 Top-1
降至约 $45\sim48\%$，非中心裁剪（\code{crop0.7off}）降至 $4\sim9\%$，
裁剪 $50\%$（\code{crop0.5}）使检索近乎失效（$<1.5\%$），
因为源图的外围背景被裁掉、逐像素对应关系被破坏。
（iii）两法的排序在光度档上完全一致，在几何档上互有胜负、差距在数个百分点的量级。
可见若只在干净图像上评测，方法之间几乎没有可分辨的差别。

\subsection{开集拒识协议}
\label{sec:openset-exp}

按第 4.5 节协议，将测试集 $200$ 条真实视频按编号切为闭集 $150$ / 开集 $50$：
图库为闭集视频的真实帧（$1200$），闭集查询（$1496$）应在库内命中，
开集查询（$504$，其源视频整体排除）应被拒绝。表~\ref{tab:openset-clean} 给出干净档读数，
表~\ref{tab:openset-deg} 给出 12 档退化下的拒识读数。

\begin{table}[t]
\centering
\caption{开集拒识（干净档）。拒识判据为对全库的 top-1 余弦，阈值在开集分布 FPR$=1\%$ 分位标定。}
\label{tab:openset-clean}
\small
\begin{tabular}{lcccccc}
\toprule
变体 & 闭集帧 T1 & 闭集帧 T5 & 视频 T1 & AUC & TPR@FPR$1\%$ & 闭/开集相似度 \\
\midrule
I-best         & 60.36 & 89.10 & 99.73 & 99.75 & 98.73 & 0.989 / 0.717 \\
r8\_norm（零训练）& 61.03 & 88.90 & 99.73 & 99.73 & 98.86 & 0.984 / 0.721 \\
\bottomrule
\end{tabular}
\end{table}

\begin{table}[t]
\centering
\caption{12 档退化下的拒识读数（AUC / TPR@FPR$1\%$）。}
\label{tab:openset-deg}
\small
\begin{tabular}{lcccc}
\toprule
\multirow{2}{*}{档位} & \multicolumn{2}{c}{I-best} & \multicolumn{2}{c}{r8\_norm（零训练）} \\
\cmidrule(lr){2-3}\cmidrule(lr){4-5}
 & AUC & TPR@1\% & AUC & TPR@1\% \\
\midrule
clean      & 99.75 & 98.73 & 99.73 & 98.86 \\
jpeg50     & 99.75 & 98.80 & 99.73 & 98.86 \\
jpeg20     & 99.74 & 98.53 & 99.73 & 98.86 \\
scale0.5   & 99.75 & 98.66 & 99.73 & 98.86 \\
crop0.5    & 55.84 & 6.15  & 51.22 & 1.94 \\
reencode   & 99.75 & 98.60 & 99.73 & 98.86 \\
\midrule
crop0.8    & 87.41 & 55.88 & 82.02 & 18.85 \\
crop0.7off & 67.04 & 7.55  & 61.53 & 14.10 \\
zoom1.15   & 98.22 & 81.28 & 97.53 & 76.74 \\
rot5       & 99.55 & 94.18 & 99.66 & 98.40 \\
rot10      & 96.87 & 64.64 & 97.71 & 76.00 \\
shift8     & 99.58 & 94.52 & 99.58 & 96.59 \\
\bottomrule
\end{tabular}
\end{table}

\textbf{结论。}（i）\textbf{干净档下拒识非常容易}：闭集与开集查询的 top-1 相似度均值相差
$0.26\sim0.27$（约 $0.99$ 与 $0.72$），AUC 达 $99.7\%$，
说明闭集查询因背景与库内源图几乎逐像素相同而被轻易识别；
此时拒识近乎平凡，同样不足以区分方法。
（ii）\textbf{退化把拒识变成一个真正的评测点}：光度档下拒识几乎不变（AUC 仍 $99.7\%$），
而几何档下显著下降，\code{crop0.5} 时 AUC 由 $99.75\%$ 降至 $55.84\%$，
TPR 仅 $6.15\%$；\code{crop0.7off} 的 AUC 虽仍有 $67\%$，TPR 却只有 $7.55\%$，
即在高召回要求下几乎不可用。
（iii）两法在拒识上的表现整体相当、互有胜负，
说明在这一维度上学习的收益同样有限。

\subsection{可微几何验证路径}
\label{sec:geo-exp}

第 5.3 节的假设指出，缺失的环节是“成对比较”。本节据此构造两阶段重排器
（廉价初筛 Top-$K$ $\rightarrow$ 候选内几何重排，第 4.3--4.4 节），
在几何七档（\code{clean} 与几何六档）下与经典 SIFT$+$RANSAC 并列比较，结果见表~\ref{tab:geo}。

\begin{table}[t]
\centering
\caption{两种几何验证路线在几何七档上的帧级 Top-1（$K=100/300$，\%）。}
\label{tab:geo}
\small
\begin{tabular}{lcc}
\toprule
阶段 / 方法 & 干净档 & 几何七档均值 \\
\midrule
初筛（融合描述子）                     & 60.00 & -- \\
可微几何重排（独立局部描述子）          & 67.00 & 66.52 \\
$+$ 与全局描述子融合（最优权重）        & \textbf{67.67} & \textbf{66.90} \\
经典 SIFT$+$RANSAC                     & -- & 67.57 \\
\bottomrule
\end{tabular}
\end{table}

\textbf{结论。}（i）\textbf{成对比较确实补上了缺失的环节}：干净档上初筛为 $60.00\%$，
加入可微几何重排后升至 $67.67\%$（$+7.7$ 个百分点）。初筛本身已是本协议下最强的一级方案，
因此这一提升并非来自更强的描述子，而是来自“逐候选做局部校验”这一形式，
与第 5.3 节的假设一致。在只保留同视频 $8$ 选 $1$ 的纯细节口径下，同一结论成立：
初筛为 $60.11\%$（零训练）$/$ $59.79\%$（学习），接入几何重排后升至 $66.25\%$，
而该口径的随机基线仅为 $12.5\%$。
（ii）\textbf{与经典路线相当}：几何七档均值 $66.90\%$ 相对 SIFT$+$RANSAC 的 $67.57\%$
相差 $0.67$ 个百分点，二者属于同一水平。逐档看，两条路线互有胜负
（条件命中率 $63\sim70\%$ 与 $62\sim72\%$），差距主要来自 \code{crop0.7off} 的初筛召回
仅为 $91.33\%$（上界受损），以及 \code{shift8}（$-6.3$ 个百分点）与 \code{clean}
（$-3.7$ 个百分点）。
（iii）第 5.8 节的参数扫描显示，这一差距不是参数未调好所致。
因此在该任务上，全可微的几何验证是一条与经典方法同等可行的路线，
其价值在于打分链路内部全程可导、无需特征检测器与随机采样。
但需要如实指出：两路局部描述子为独立训练，几何损失的梯度并不回传全局编码器；
把几何一致性并入训练的尝试（K/K2 轮的辅助损失与独立网络）均未带来增益，见第 5.8 节。

\chead{扩展退化档上的配对检验。}上述七档只覆盖几何退化，而模糊、低分辨率与极端压缩
同样常见。为此我们在 $1000$ 条查询的分层子集上补测 $11$ 档（含 \code{blur2}、
\code{blur4}、\code{lowres64}、\code{jpeg5}），并按源视频聚类做 bootstrap 与 McNemar
配对检验。本文两阶段均值 $68.35\%$，SIFT$+$RANSAC 为 $60.68\%$，差值 $+7.67$ 个百分点
（$95\%$ CI $[+6.0,+9.4]$），$11$ 档全部胜出，$9/11$ 档 $p<0.05$。
优势集中在 \code{blur4}（$+21.0$）、\code{blur2}（$+15.1$）、\code{lowres64}（$+14.2$）
与 \code{crop0.7off}（$+11.9$）；\code{clean}（$+0.7$，$p=0.52$）与
\code{shift8}（$+1.3$，$p=0.21$）不显著，即干净档与纯平移两法持平。
该子集与全量协议（图库 $1600$ / 查询 $8000$）不可混用，故单列于此。
还需要说明，在真实 H.264 二次转码的补充对照中两法打平（CRF23 与 CRF32 的差值
分别为 $-0.9/+0.2$ 个百分点），因此 $+7.67$ 不能外推为“真实传播链路也领先”。

\subsection{真帧级定位：稠密帧库与相似度命中区间}
\label{sec:dense-exp}

\chead{为什么需要这一实验。}前面各节的“帧级”判据是：查询取自某条视频的某个帧号，
命中判为图库中\textbf{同一帧号}的那张真实帧。由于库中每条视频只取 $8$ 帧（相隔约两秒），
这一判据实际衡量的是“在相隔两秒的 $8$ 张里挑对一张”，而非“在一段连续视频中定位伪造帧来自哪一时刻”。
对真实取证而言后者才是需要的能力——调查者要拿一张伪造帧回到原始素材里定位具体时刻。
本节因此重建协议：图库由连续帧构成，命中按\textbf{帧相似度}而非帧号判定。

\chead{构造。}沿用第 3.2 节的测试划分（\code{original} 视频按名排序后取第 $400$–$599$ 条），
从中均匀取 $30$ 条，每条自第 $100$ 帧起连续抽 $100$ 帧，
逐帧用 YuNet 检测最大人脸并外扩 $1.8$ 倍裁为 $256\times256$，构成 $3000$ 帧图库。
逐帧检测是必要的：同一条视频内人脸框随姿态移动（$x$ 坐标极差均值约 $105$ 像素），
沿用整段一个框会让“同帧”的裁切错位。查询取这 $30$ 条视频对应的
Deepfakes$/$Face2Face$/$FaceSwap 伪造帧各 $6$ 帧，共 $540$ 张。
不加退化即 \code{clean} 档，退化档与第 5.4 节相同且只作用于查询。

\chead{命中判据。}若以“帧号相等”为命中，会系统性低估：换脸方法改变帧率与节奏，
伪造第 $t$ 帧的背景未必取自真实第 $t$ 帧（实测偏移中位为 $0$，绝对偏移中位 $2$、$p90$ 达 $12$）；
静止段内相邻帧又几乎相同，帧号差一两帧在视觉上不可区分。因此把真值定义为\textbf{相似度区间}：
对查询 $q$，令 $d(q,g)=1-\cos(\mathrm{r8}(q),\mathrm{r8}(g))$（$8\times8$ 逐图标准化缩略图，第 5.2 节），
\begin{equation}
\mathrm{GT}(q)=\{\,g:\ g\ \text{与}\ q\ \text{同视频},\ d(q,g)\le d_{\min}(q)+\delta\,\},
\end{equation}
$d_{\min}$ 为同视频内最近距离，$\delta$ 取同视频相邻帧距离分布的 $p50$、$p90$、$p95$
（$0.0008$、$0.0060$、$0.0119$），依次对应“几乎只认同帧”“容许约 $4$ 帧”“容许约 $23$ 帧”；
$|{\rm GT}|$ 的中位数分别为 $3.8$、$23.2$、$37.7$，而每视频有 $100$ 帧，
故\emph{同视频随机基线}为 $3.8\%$、$23.2\%$、$37.7\%$。需要强调，
$\mathrm{r8}$ 只用于\textbf{定义真值}，不作为检索方法报告，否则等于用判据本身作答。

\chead{指标与对照。}报告 $T1_{\text{video}}$（top-$1$ 是否落在正确视频）、
$T1_{\text{loc}}$（是否落进 $\mathrm{GT}$）、$T5_{\text{loc}}$、条件定位率（视频命中条件下的 $T1_{\text{loc}}$）、
top-$20$ 召回，以及定位误差 $|\Delta|$（预测帧号与最近真值帧号之差）的中位数与 $\le 2$ 帧比例。
三种方法共用同一图库、查询与初筛：全局单路（I-best 嵌入余弦，即两阶段的第一阶段）、
本文两阶段（$K=20$ 几何重排）、SIFT$+$RANSAC（与本文共用同一初筛）。
表~\ref{tab:dense} 给出五档读数；$\mathrm{r8}$ 因参与定义真值而不列入。

\begin{table}[t]
\centering
\caption{稠密帧库上的真帧级定位（图库 $3000$ / 查询 $540$，$K=20$，\%）。
$T1_{\text{loc}}$ 三列依次对应 $\delta=p50/p90/p95$；$T5_{\text{loc}}$ 与 $\le2$ 帧取 $\delta=p50$。
随机基线（同视频）依次为 $3.8/23.2/37.7$。}
\label{tab:dense}
\small
\begin{tabular}{llcccccc}
\toprule
档位 & 方法 & $T1_{\text{video}}$ & $T1_{\text{loc}}$ p50 & p90 & p95 & $T5_{\text{loc}}$ & $\le2$ 帧 \\
\midrule
\multirow{3}{*}{\code{clean}}
 & 全局单路      & 100.0 & \textbf{85.2} & \textbf{98.9} & \textbf{99.8} & \textbf{97.4} & \textbf{89.3} \\
 & 本文两阶段    & 100.0 & 65.9 & 94.1 & 98.5 & 91.7 & 82.2 \\
 & SIFT$+$RANSAC & 100.0 & 48.5 & 90.4 & 96.5 & 86.5 & 76.3 \\
\midrule
\multirow{3}{*}{\code{crop0.8}}
 & 全局单路      & 97.6 & 5.6 & 19.1 & 32.6 & 17.6 & 16.9 \\
 & 本文两阶段    & 97.8 & \textbf{19.4} & 51.3 & 68.1 & \textbf{35.2} & \textbf{44.3} \\
 & SIFT$+$RANSAC & 97.8 & 17.8 & \textbf{53.7} & \textbf{72.2} & 33.7 & 40.9 \\
\midrule
\multirow{3}{*}{\code{crop0.7off}}
 & 全局单路      & 62.2 & 3.5 & 11.3 & 20.9 & 9.6 & 12.0 \\
 & 本文两阶段    & \textbf{74.1} & \textbf{12.4} & \textbf{40.2} & \textbf{51.9} & \textbf{28.9} & \textbf{33.3} \\
 & SIFT$+$RANSAC & \textbf{74.1} & 10.9 & 30.6 & 45.4 & 20.4 & 25.2 \\
\midrule
\multirow{3}{*}{\code{blur4}}
 & 全局单路      & 100.0 & \textbf{83.3} & \textbf{98.9} & \textbf{100.0} & \textbf{97.2} & \textbf{88.0} \\
 & 本文两阶段    & 100.0 & 51.1 & 92.0 & 97.2 & 87.6 & 76.3 \\
 & SIFT$+$RANSAC & 100.0 & 39.1 & 86.1 & 95.9 & 78.1 & 65.7 \\
\midrule
\multirow{3}{*}{\code{jpeg5}}
 & 全局单路      & 100.0 & \textbf{54.8} & \textbf{93.1} & \textbf{98.3} & 83.0 & 71.7 \\
 & 本文两阶段    & 100.0 & 51.9 & 90.6 & 98.0 & \textbf{88.5} & \textbf{73.7} \\
 & SIFT$+$RANSAC & 100.0 & 39.8 & 88.0 & 95.9 & 80.4 & 67.0 \\
\bottomrule
\end{tabular}
\end{table}

\textbf{结论。}（i）\textbf{真帧级定位可行，但难度取决于判据}：$p90$ 区间下三种方法在
\code{clean} 档分别为 $98.9\%$、$94.1\%$、$90.4\%$，远高于同视频随机基线 $23.2\%$；
区间收紧到 $p50$ 后仍有 $85.2\%/65.9\%/48.5\%$，对随机基线 $3.8\%$。
（ii）\textbf{干净与光度档由全局外观主导}：\code{clean}、\code{blur4}、\code{jpeg5} 上全局单路的
$T1_{\text{loc}}$ 全面最高（$p50$ 为 $85.2/83.3/54.8$），两种几何重排都更低。
这与第 5.6 节“几何重排提升帧级”的方向相反，原因在判据不同：旧协议要求区分相隔约两秒的 $8$ 帧，
而稠密库中最近邻常只差一两帧，全局内积已足够，局部几何校验反而把分数推向纹理更相似的错误近邻。
（iii）\textbf{几何退化档才是几何重排的用武之地}：\code{crop0.8} 上全局单路跌至
$5.6\%/19.1\%/32.6\%$，本文两阶段为 $19.4\%/51.3\%/68.1\%$（$p50$ 上 $+13.8$ 个百分点）；
\code{crop0.7off} 上全局单路仅 $3.5\%/11.3\%/20.9\%$，本文两阶段为 $12.4\%/40.2\%/51.9\%$，
视频级命中由 $62.2\%$ 升至 $74.1\%$。
（iv）\textbf{与 SIFT 互有胜负}：\code{crop0.8} 上 SIFT 在 $p90/p95$ 略高（$+2.4/+4.1$），
\code{crop0.7off} 上本文更高（$+6.5$ 个百分点，$p95$），其余档位差别在数个点内，
与第 5.6 节“两条几何路线同一水平”的判断一致。
（v）\textbf{需要修正的设计点}：\code{blur4} 上本文两阶段仅 $51.1\%$，明显低于全局单路的 $83.3\%$，
而该档门控回退触发率达 $39.6\%$。这说明第 4.5 节的门控与几何回退阈值是在旧协议上标定的，
不能直接迁移到稠密协议，应按判据重标定；本文如实报告该负增益，不以档均值掩盖。
此外，初筛 top-$20$ 召回在 \code{clean}/\code{crop0.8}/\code{crop0.7off} 为
$100\%/84.4\%/68.0\%$，即 \code{crop0.7off} 的帧级定位率上界已被初筛封定。

\subsection{消融实验}
\label{sec:ablation}

\chead{（a）几何模块参数敏感性。}
我们对软匹配温度 $\tau$ 与软内点尺度 $\sigma$ 做 $4\times4$ 网格扫描（每个组合重算候选内
几何分数），几何七档均值见表~\ref{tab:param}。

\begin{table}[t]
\centering
\caption{$\tau\times\sigma$ 网格下的几何七档帧级 Top-1 均值（纯几何 / 融合）。}
\label{tab:param}
\small
\begin{tabular}{lcccc}
\toprule
$\tau \backslash \sigma$ & 0.02 & 0.04 & 0.06 & 0.08 \\
\midrule
0.01 & 61.05 / 61.67 & 64.57 / 64.76 & 65.76 / 65.90 & 66.14 / 66.38 \\
0.03 & 62.71 / 63.05 & 64.95 / 65.38 & 66.05 / 66.33 & \textbf{66.52 / 66.90} \\
0.05 & 62.90 / 63.14 & 65.38 / 65.62 & 66.10 / 66.52 & 66.52 / 66.76 \\
0.10 & 58.57 / 59.14 & 63.52 / 63.62 & 64.81 / 65.00 & 65.19 / 65.24 \\
\bottomrule
\end{tabular}
\end{table}

\textbf{结论。}$\sigma$ \textbf{越大越好}，在四个 $\tau$ 上均以 $\sigma=0.08$ 最优；
$\tau$ 在 $0.01\sim0.05$ 区间近乎无差，$\tau=0.1$ 略差。最优组合
（$\tau=0.03,\sigma=0.08$）给出 $66.52\%$（纯几何）$/$ $66.90\%$（融合），
与经典方法相差约 $1$ 个百分点；这一差距在网格内无法消除。
这一结果说明，可微几何验证的精度上限受限于初筛召回与局部描述子本身，
而不是软匹配参数。

\chead{（b）几何一致性作为训练目标。}
我们还尝试把几何一致性分别用于训练的三个位置，结果见表~\ref{tab:geotrain}。

\begin{table}[t]
\centering
\caption{几何一致性目标在训练中的三种尝试及其结果。}
\label{tab:geotrain}
\small
\begin{tabular}{p{3.4cm}p{4.4cm}p{6.2cm}}
\toprule
位置 & 做法 & 结果 \\
\midrule
辅助损失（共享主干） & 挂到全局编码器（K / K2 轮） & 帧 T1 58.71 / 57.96，低于 I-best 59.68 \\
独立网络（自 InfoNCE 微调） & 铰链损失 $\mathrm{relu}(m+s_n-s_p)$，$m=0.3$ & 起步 $s_p\approx0.81,\ s_n\approx0.04$，分离 $\approx0.77\gg m$，铰链饱和、梯度为零 \\
独立网络（随机初始化） & 直接以几何目标训练 & $s_p\approx s_n\approx0.95$，退化解塌陷 \\
\bottomrule
\end{tabular}
\end{table}

\textbf{结论。}无论作为辅助损失（K/K2）还是独立网络的唯一目标，
几何一致性都未能带来检索增益：辅助损失两轮均低于 I-best；
独立训练时，由于正负对之间的几何分数分离度远大于铰链边界，目标直接饱和（梯度为零），
换用排序损失后分数不升反降，而随机初始化则收敛到平凡解。
这三点一致说明：在该任务中，几何一致性适合作为\textbf{检索期的打分器}（第 5.6 节），
而不适合作为\textbf{训练期的损失}。

\subsection{跨数据集：Celeb-DF v2 上的真实换脸溯源}
\label{sec:cross-exp}

前述实验全部建立在 FF++\_C23 的同域换脸协议上（查询与真值来自同一条视频，背景近乎逐像素相同），
因此无法回答一个自然的问题：换一个未见过的数据域与换脸实现，方法是否仍然成立？
本节在两个补充集合上检验：真实换脸视频语料 Celeb-DF v2（本小节），
以及图像级的区域拼接伪造（第 5.10 节）。

\paragraph{真值恢复。}Celeb-DF v2~\cite{li2020celebdf} 的伪造文件命名为 \code{idX\_idY\_000k.mp4}。
我们确认其语义为：\code{idX} 是背景（原视频）的身份、\code{idY} 是被换上去的人脸身份、
\code{k} 是该身份真实视频列表中的序号，故伪造帧的背景源恰为 \code{Celeb-real/idX\_000k.mp4}。
依据有三：全部 $5639$ 条伪造都满足对应的真实视频存在；抽查样本在背景（边框）区域的逐像素差
对声称的源为 $1.8$–$2.5$，对其余真实视频均在 $30$ 以上；两者时长逐帧一致。
由此可以恢复\textbf{精确的“伪造帧 $\rightarrow$ 背景源视频”真值}，
而不必以 real$\rightarrow$real 近似替代。

\paragraph{协议。}图库为全部 $590$ 条真实视频各取一帧（$t=0.5$ 时长），可用 $589$ 条；
查询取官方测试清单中的 $340$ 条伪造视频、各取两帧（$t=0.25$ 与 $0.75$ 时长），共 $680$ 条，
来自 $213$ 个不同背景源；帧统一缩放至 $256\times256$。
查询帧与图库帧的时刻\textbf{刻意错开}，以免“同一时刻、同一背景”构成平凡解。
退化档与第 5.4 节相同且只作用于查询，两种重排器共用同一全局初筛（$K=20$）。
随机猜中率为 $1/589=0.17\%$。表~\ref{tab:celebdf} 给出五档读数。

\begin{table}[t]
\centering
\caption{Celeb-DF v2 跨数据集溯源（图库 $589$ / 查询 $680$，命中=背景源视频，Top-1，\%）。
两种重排器共用全局初筛 top-20；退化只作用于查询。}
\label{tab:celebdf}
\small
\begin{tabular}{lccccc}
\toprule
档位 & $8\times8$ 缩略图 & 全局单路 & \textbf{本文（两阶段）} & SIFT$+$RANSAC & 差值 \\
\midrule
clean      & 95.15 & 97.79 & \textbf{98.82} & 98.24 & $+0.59$ \\
crop0.8    & 82.06 & 89.85 & \textbf{98.09} & 96.62 & $+1.47$ \\
crop0.7off & 15.88 & 38.24 & \textbf{77.06} & 75.29 & $+1.76$ \\
blur4      & 95.15 & 97.79 & \textbf{97.79} & 93.53 & $\mathbf{+4.26}$ \\
jpeg5      & 95.15 & \textbf{95.44} & 93.53 & 93.82 & $-0.29$ \\
\midrule
均值       & 76.68 & 83.82 & \textbf{93.06} & 91.50 & $+1.56$ \\
\bottomrule
\end{tabular}
\end{table}

\paragraph{结论。}（i）\textbf{跨数据集不掉点}：干净档 $98.82\%$，高于零训练缩略图 $3.7$ 个百分点、
高于全局单路 $1.0$ 个百分点，且两种重排器都超过 $98\%$。
Celeb-DF 的伪造来自真实网络素材并经 H.264 转码，这一读数说明“以背景为来源指纹”的假设
在真实换脸域同样成立，而不只是在自建合成集合上成立。
（ii）\textbf{优势集中在几何裁剪，且相对自身初筛的增益远大于与 SIFT 的差距}：
\code{crop0.7off} 相对全局单路提升 $38.8$ 个百分点、相对零训练缩略图提升 $61.2$ 个百分点，
但与 SIFT 仅差 $1.76$ 个百分点。原因是该档共用初筛的召回已塌至 $77.35\%$，
两种重排器的候选内条件 Top-1 分别为 $98.50\%$ 与 $97.34\%$，差距被初筛封顶。
（iii）\code{blur4} 领先 SIFT $4.26$ 个百分点：该档 SIFT 候选内的中位好匹配数为 $0.4$、
中位内点数为 $0$，可微描述子仍维持 $98.37\%$ 的条件 Top-1，与第 5.10 节的拼接集结论一致。
（iv）\code{jpeg5} 是本协议下唯一未领先的档位（$-0.29$ 个百分点，本文 $93.53\%$ 低于全局单路 $95.44\%$）：
该档背景在压缩后仍近乎无损（零训练缩略图仍有 $95.15\%$），几何重排没有可补的信息，
其随机噪声反而略损排序。这与表~\ref{tab:degrade} 中同域 \code{jpeg5} 的增益方向不同，
说明光度档的胜负取决于压缩是否真的破坏了布局，而不能笼统断言“光度档全胜”。

\paragraph{对照：换脸本身是否增加难度。}在同协议的 real$\rightarrow$real 对照上
（同一 $213$ 个背景源、同样的两个时刻，仅查询帧未换脸，共 $426$ 条；图库与判定完全一致），
干净档与 \code{crop0.7off} 分别为 $98.36\%$ 与 $78.64\%$，
与伪造帧的 $98.82\%$ 与 $77.06\%$ 相差在 $0.5$–$1.6$ 个百分点内。
即：换脸几乎不损失可溯源性，本任务的难度来自“同视频不同帧 + 退化”，而不是“脸被换掉”。
这既支持了本文以背景为指纹的建模选择，也意味着该表不能用来论证方法对伪造痕迹本身敏感。

\subsection{跨伪造类型：区域拼接压力测试}
\label{sec:splice-exp}

上一个小节换了数据域与换脸实现，但伪造类型仍是自编码器式换脸。
为检验该假设对别的篡改方式是否成立，我们用 CelebAMask-HQ~\cite{lee2019maskgan}
构造一个\textbf{补充压力测试}：图库为 $929$ 张干净的真实人脸图（$224\times224$，每张图对应唯一身份），
查询为 $929$ 张拼接图——保留背景图 $A$，按掩码把他人的脸区（脸皮与五官，按 bbox 对齐、
边界羽化）贴入，命中真值仍为 $A$。该集合改变了伪造类型（splice），
但\textbf{不能替代}视频级换脸实验：它是图像级、每源仅一张图库图，且拼接集为自建而非标准 benchmark。

评测沿用第 4.4 节的两阶段管线（本文最终配置：$64$ 维局部描述子与 ResNet-18 第二路融合、
门控与几何回退同第 4.4 节实现，$K=20$），并让 SIFT$+$RANSAC~\cite{lowe2004sift,fischler1981ransac}
与本文方法\textbf{共用同一个全局初筛}，唯一差异是重排器。表~\ref{tab:splice} 给出五档读数。

\begin{table}[t]
\centering
\caption{CelebA-HQ 拼接压力测试（图库 $929$ / 查询 $929$，图像级，帧级 Top-1，\%）。
两种重排器共用全局初筛 top-20；退化只作用于查询。}
\label{tab:splice}
\small
\begin{tabular}{lcccc}
\toprule
档位 & $8\times8$ 缩略图 & 全局单路 & \textbf{本文（两阶段）} & SIFT$+$RANSAC \\
\midrule
clean      & 89.77 & 88.81 & \textbf{96.99} & 96.99 \\
crop0.8    & 11.30 & 11.73 & \textbf{39.18} & 38.97 \\
crop0.7off & 6.03  & 9.47  & \textbf{27.23} & 27.02 \\
blur4      & 89.45 & 88.05 & \textbf{96.88} & 79.55 \\
jpeg5      & 90.20 & 85.68 & \textbf{96.02} & 89.88 \\
\midrule
均值       & 57.35 & 56.75 & \textbf{71.26} & 66.48 \\
\bottomrule
\end{tabular}
\end{table}

\textbf{结论。}（i）\textbf{跨域与跨伪造类型均未使系统失效}：干净档 $96.99\%$，
高于零训练缩略图 $7.2$ 个百分点、高于全局单路 $8.2$ 个百分点；且只要真值进入候选集，
重排后的条件 Top-1 在五档均不低于 $99.5\%$，说明拼接把“同源不同脸”的差异放大后，
几何重排反而比在同域帧级任务中更决断。
（ii）\textbf{光度档的优势复现}：\code{blur4} 领先 SIFT $17.33$ 个百分点、
\code{jpeg5} 领先 $6.14$ 个百分点。原因与第 5.4 节一致——模糊抹掉了 SIFT 依赖的梯度角点
（该档 SIFT 候选内的中位好匹配数为 $0$），而可微描述子仍能建立对应。
（iii）\textbf{几何裁剪档与 SIFT 打平}（$+0.22$ 与 $+0.22$ 个百分点，约 $2/929$ 条）：
两侧同时被共用的全局初筛卡住，候选召回仅 $39.50\%$ 与 $28.31\%$，
而 SIFT 的条件 Top-1 仍有 $98.64\%$ 与 $95.44\%$——瓶颈在召回而非重排，
这与第 5.4 节对 \code{crop0.7off} 的诊断同构。因此不能据此声称在跨域裁剪下领先。
（iv）本表的\textbf{绝对数值不可与表~\ref{tab:degrade} 并列}：每个来源只有一张图库图
（没有同视频多帧的多次机会），拼接区域平均占 $36.4\%$ 的像素，且为图像级而非帧级任务，
\code{crop0.8} 的 $39.18\%$ 因此明显低于同域协议的 $69.8\%$，这是协议差异而非方法退化。
综上，本小节只提供“同一结构在拼接伪造下复现”的补充证据；
视频级的真实换脸跨数据集验证已由第 5.9 节给出，DFDC 与修复类伪造仍未覆盖。
